\section{Fee Sharing Protocol}

\label{sec:feesharing}
%%==========================================================================

\subsection{Motivation}

Lux appchains currently operate under a weak economic coupling to the Primary Network:
validators must hold Primary Network stake (2,000 LUX minimum) but appchains pay no
ongoing fees to the Primary Network for security. This creates a free-rider problem:
large appchains derive substantial security from the global LUX staking pool without
contributing to its maintenance.

We propose an \emph{Appchain Security Fee} (ASF) mechanism that transfers a fraction of
appchain fee revenue to Primary Network validators.

\subsection{Two-Tier Fee Structure}

\paragraph{Tier 1: Execution Fees.}
Fees collected by appchain validators for transaction execution and block production.
These remain entirely with the appchain validator set:
\[
  F_{\text{exec}} = \sum_{\text{txn}} \text{GasUsed} \cdot \text{GasPrice}.
\]

\paragraph{Tier 2: Security Surcharge.}
An additional per-transaction surcharge $\phi_s$ paid to Primary Network validators.
This is denominated in LUX and burned/redistributed:
\[
  F_{\text{security}}(s) = \phi_s \cdot \text{TxCount}(s, \text{epoch}).
\]

\subsection{Security Surcharge Formula}

The security surcharge $\phi_s$ is computed as:
\begin{equation}
  \phi_s = \phi_0 \cdot \left(\frac{\Pi(s)}{\bar{\Pi}}\right)^\gamma,
  \label{eq:surcharge}
\end{equation}

where $\phi_0$ is the base surcharge (0.001 LUX per transaction), $\Pi(s)$ is the economic
security of appchain $s$ (as defined in the Lux Validator Economics companion paper), $\bar{\Pi}$ is the median
security across all appchains, and $\gamma \in (0,1)$ is a concavity parameter (default 0.5)
ensuring that high-security appchains pay more but not in proportion.

\subsection{Revenue Distribution}

Security surcharge revenue is distributed as follows:

\begin{itemize}[leftmargin=1.5em]
  \item 60\%: Primary Network validators (proportional to stake weight);
  \item 20\%: LUX buyback-and-burn (deflationary);
  \item 15\%: LUX development fund (governed by LUX DAO);
  \item 5\%: LRP insurance fund (see the Lux Restaking Protocol companion paper).
\end{itemize}

\subsection{Retrospective Analysis}

We simulate the ASF mechanism on 2.5 years of historical Lux mainnet data (August 2023--
2025). Assuming $\phi_0 = 0.001$ LUX and $\gamma = 0.5$:

\begin{table}[H]
\centering
\caption{Simulated ASF Revenue (Historical Calibration)}
\label{tab:ssf}
\begin{tabular}{lrrr}
\toprule
\textbf{Period} & \textbf{Appchain Tx Volume} & \textbf{Simulated ASF} & \textbf{Primary Net Yield Impact} \\
\midrule
Q3 2023 & 42M & \$420K & $+$0.08\% annualized \\
Q4 2023 & 78M & \$780K & $+$0.14\% annualized \\
Q1 2024 & 134M & \$1.34M & $+$0.24\% annualized \\
Q2 2024 & 215M & \$2.15M & $+$0.39\% annualized \\
Q3 2024 & 312M & \$3.12M & $+$0.56\% annualized \\
Q4 2024 & 489M & \$4.89M & $+$0.87\% annualized \\
Q1 2025 & 612M & \$6.12M & $+$1.09\% annualized \\
Q2 2025 & 834M & \$8.34M & $+$1.48\% annualized \\
\bottomrule
\end{tabular}
\end{table}

By Q2 2025, ASF would contribute approximately 1.48\% to the annualized Primary Network
validator yield, a meaningful incentive for continued Primary Network participation as
baseline LUX inflation decreases over time.

\subsection{Positive Feedback Dynamics}

The ASF mechanism creates a positive feedback loop:
\[
  \text{More appchains} \to \text{More ASF revenue} \to \text{Higher primary yield} \to
  \text{More validators} \to \text{More security} \to \text{Easier appchain bootstrapping}.
\]

We model this as a differential equation for Primary Network stake $W(t)$:
\begin{equation}
  \frac{dW}{dt} = \lambda_W \cdot \left[y_{\text{primary}}(W) + y_{\text{ASF}}(t) - y_{\text{market}}\right] \cdot W,
  \label{eq:stakeflow}
\end{equation}

where $y_{\text{primary}}(W)$ is the baseline staking yield (decreasing in $W$),
$y_{\text{ASF}}(t)$ is the ASF contribution (increasing in appchain volume), and
$y_{\text{market}}$ is the market return on capital. The equilibrium $W^*$ satisfies
$y_{\text{primary}}(W^*) + y_{\text{ASF}}(t^*) = y_{\text{market}}$. As appchain volume
grows, $y_{\text{ASF}}(t)$ increases, pushing $W^*$ upward.

%%==========================================================================
